% 3.3 — Architecture d'ensemble
\section{Architecture d'ensemble}
\label{sec:architecture}

\subsection{Le flux d'une commande}
La figure~\ref{fig:flux} suit une commande de l'agent jusqu'au reçu. Cette section présente
l'architecture par ses idées ; l'annexe~\ref{ann:architecture} la montre en entier, module par
module.

\begin{figure}[htbp]
\centering
% Figure 1 — Le flux d'ensemble : de l'agent au reçu.
\begin{tikzpicture}[node distance=5mm and 6mm]
  % --- Rangée du haut : agent -> portes -> moteur
  \node[douce,text width=26mm] (agent) at (0,0) {Agent de code IA\sous{Claude Code, Codex}};
  \node[boite,text width=30mm] (edit) at (4.6,0) {Porte fichier\sous{Edit, Write}};
  \node[boite,text width=30mm,above=2mm of edit] (hook) {\textit{Hook} PreToolUse\sous{Bash, PowerShell}};
  \node[boite,text width=30mm,below=2mm of edit] (shell) {Garde du shell\sous{piège bash, \texttt{check}}};
  \node[zone,fit=(hook)(edit)(shell),label={[etiquette]above:portes}] (portes) {};
  \draw[fleche] (agent.east) -- (portes.west |- agent.east);

  \node[accent,text width=36mm] (sim) at (9.9,0) {Simulation\sous{FS virtuel (\texttt{simulate.py})}};
  \node[accent,text width=36mm,above=3mm of sim] (parse) {Analyse syntaxique\sous{POSIX \textbar{} PowerShell (C3)}};
  \node[accent,text width=36mm,below=3mm of sim] (eff) {Effets $\rightarrow$ impact\sous{détruit, crée, distant, inconnu}};
  \node[zone,draw=wepwnfonce,fit=(parse)(sim)(eff),label={[etiquette,text=wepwnfonce]above:moteur}] (moteur) {};
  \draw[fleche] (parse) -- (sim);
  \draw[fleche] (sim) -- (eff);
  \draw[fleche] (portes.east |- sim.west) -- (moteur.west |- sim.west);

  % --- Politique et décideur
  \node[boite,text width=36mm] (pol) at (9.9,-3.4) {\texttt{verdict} + \texttt{policy}\sous{effet $\times$ surface}};
  \node[boite,draw=encre,line width=1.1pt,text width=30mm] (dec) at (4.6,-3.4) {\textbf{\texttt{decide.py}}\sous{le décideur unique}};
  \draw[fleche] (moteur.south) -- (pol.north);
  \draw[fleche] (pol.west) -- (dec.east);

  % --- Décisions (colonnes fixes)
  \foreach \x/\n/\c/\t in {0/o1/dALLOW/ALLOW,
                          2.6/o2/dREV/{REVERSIBLE\\[-2pt]{\tiny ou DRY\_RUN}},
                          5.2/o3/dSUP/SUPERVISED,
                          7.8/o4/dCTX/{CONTEXT\\[-2pt]MISMATCH},
                          10.4/o5/dASK/{ESCALATE\\[-2pt]{\tiny ask ou deny}}}
    \node[decision,fill=\c,text width=21mm] (\n) at (\x,-5.4) {\t};
  \draw[fleche] (dec.south) -- ++(0,-0.45) coordinate (b);
  \draw[encre!70,line width=0.7pt] (o1.north |- b) -- (o5.north |- b);
  \foreach \n in {o1,o2,o3,o4,o5} \draw[fleche] (\n.north |- b) -- (\n.north);

  % --- Conséquences
  \node[douce,text width=21mm] (rec) at (2.6,-7.0) {Point de\\récupération\sous{\texttt{undo <id>}}};
  \node[douce,text width=21mm] (net) at (5.2,-7.0) {Filet\\d'exécution\sous{ptrace / WinFsp}};
  \node[douce,text width=21mm] (per) at (9.1,-7.0) {La personne\\tranche\sous{ou refus net}};
  \draw[fleche,densely dashed] (o2) -- (rec);
  \draw[fleche,densely dashed] (o3) -- (net);
  \draw[fleche,densely dashed] (o4.south) -- (o4.south |- per.north);
  \draw[fleche,densely dashed] (o5.south) -- (o5.south |- per.north);

  % --- Reçu
  \node[boite,draw=wepwnfonce,line width=0.9pt,minimum width=128mm] (recu) at (5.2,-8.7)
       {Reçu dans le registre chaîné (SHA-256) --- \emph{chaque} décision en laisse un};
  \foreach \n in {o1,rec,net} \draw[fleche] (\n.south) -- (\n.south |- recu.north);
  \draw[fleche] (per.south) -- (per.south |- recu.north);
\end{tikzpicture}

\caption{Le flux d'ensemble : de la commande de l'agent au reçu}
\label{fig:flux}
\end{figure}

\begin{enumerate}
  \item \textbf{La porte} reçoit l'action : le \textit{hook} PreToolUse de Claude Code (outils
    \cmd{Bash} et \cmd{PowerShell}), la porte fichier (\cmd{Edit}, \cmd{Write}), le
    \textit{hook} de Codex ou de Cursor, le piège du shell pour les commandes tapées dans un
    terminal, ou \cmd{demo\_cli check}. Chaque porte ne fait que traduire : elle passe la
    commande à l'orchestrateur et traduit la décision dans ce que son agent comprend.
  \item \textbf{L'orchestrateur} (\fichier{guard.py}, \cmd{Guard.evaluate}) enchaîne des
    étapes nommées : refus d'une configuration cassée ; prédiction du moteur ; refus de ce
    qui touche le plan de contrôle ; choix du chemin (moteur, ligne non lue, ancien chemin) ;
    contexte ; décision ; reçu.
  \item \textbf{Le moteur} analyse la ligne, la simule sur une vue virtuelle du système de
    fichiers et produit un \emph{impact} : ce qui est détruit, créé, atteint à distance, non
    compris, non couvert.
  \item \textbf{La politique} (\fichier{policy.py}) traduit chaque effet distant en
    décision, chaque étiquette en surface, et chaque surface en canal « demander » ou
    « refuser ».
  \item \textbf{Le décideur} (\fichier{decide.py}), fonction pure sans entrée-sortie, rend la
    décision. Il ne capture rien lui-même : il reçoit seulement le fait qu'une capture a
    \emph{réussi}. C'est ce qui l'empêche de mentir.
  \item \textbf{Les conséquences} : une capture (REVERSIBLE), un filet (SUPERVISED), une
    question à la personne ou un refus (ESCALATE) --- et, dans tous les cas, un reçu.
\end{enumerate}

\subsection{Ce que l'outil voit}
\label{sec:vision}
La figure~\ref{fig:vision} répond à la question qu'un jury pose en premier : \emph{que
voit réellement l'outil, et où s'arrête-t-il ?} Elle croise six couches --- lire,
récupérer, surveiller, refuser, demander, tracer --- avec trois zones où une commande peut
agir.

\begin{figure}[htbp]
\centering
% Figure 2 — Ce que l'outil voit : chaque couche, et ce qu'elle couvre (ou pas),
% selon l'endroit où la commande agit.
\begin{tikzpicture}[
    cell/.style={draw=white,line width=1.2pt,minimum width=38mm,minimum height=12.5mm,
                 text width=35mm,align=center,font=\sffamily\scriptsize,inner sep=2pt},
    oui/.style={cell,fill=wepwnclair},
    part/.style={cell,fill=dASK!18},
    non/.style={cell,fill=grisclair,text=gris,
                postaction={pattern=north east lines,pattern color=trait}},
    ligne/.style={minimum width=34mm,minimum height=12.5mm,text width=32mm,
                  align=left,font=\sffamily\small,inner sep=2pt},
    tete/.style={minimum width=38mm,font=\sffamily\small\bfseries,align=center,text width=36mm},
  ]
  \def\W{3.85}\def\H{1.3}
  % En-têtes de colonnes
  \node[tete] at (1*\W,0.2) {Dans le projet};
  \node[tete] at (2*\W,0.2) {Sur la machine,\\hors du projet};
  \node[tete] at (3*\W,0.2) {Surfaces distantes\\{\scriptsize\mdseries base, cloud, git, SaaS}};

  % Lignes : (y, titre, sous-titre)
  \foreach \y/\t/\s in {1/Lire avant/le moteur (statique),
                        2/Récupérer/{instantané, point de contrôle},
                        3/Surveiller/{filet : ptrace ou WinFsp},
                        4/Refuser/plan de contrôle,
                        5/Demander/la personne (ask),
                        6/Tracer/les reçus}
    \node[ligne] at (0.15,-\y*\H) {\textbf{\t}\\{\scriptsize\color{gris}\s}};

  % Lire
  \node[oui]  at (1*\W,-1*\H) {prédit ce qui est détruit,\\créé, déplacé};
  \node[oui]  at (2*\W,-1*\H) {prédit la destruction\\$\rightarrow$ ESCALATE};
  \node[oui]  at (3*\W,-1*\H) {reconnaît la surface\\et la sorte d'effet};
  % Récupérer
  \node[oui]  at (1*\W,-2*\H) {copie exacte avant,\\tout ou rien};
  \node[non]  at (2*\W,-2*\H) {aucune copie\\(V17 : jamais prétendue)};
  \node[non]  at (3*\W,-2*\H) {aucune copie locale ;\\indice chez le fournisseur};
  % Surveiller
  \node[oui]  at (1*\W,-3*\H) {sauvegarde chaque fichier\\au moment où il est détruit};
  \node[part] at (2*\W,-3*\H) {enregistré, non sauvegardé\\(V18, filet Linux)};
  \node[non]  at (3*\W,-3*\H) {invisible au filet\\local};
  % Refuser
  \node[oui]  at (1*\W,-4*\H) {\texttt{.demo\_cli/}, \texttt{.demo\_cli.toml},\\\textit{hooks} du projet};
  \node[part] at (2*\W,-4*\H) {chemins listés seulement :\\\texttt{\textasciitilde/.bashrc}, \textit{hooks}, tâche};
  \node[non]  at (3*\W,-4*\H) {sans objet};
  % Demander
  \node[part] at (1*\W,-5*\H) {seulement si rien\\n'a pu être capturé};
  \node[part] at (2*\W,-5*\H) {destruction hors projet\\: ask};
  \node[oui]  at (3*\W,-5*\H) {ask ou deny\\selon la surface};
  % Tracer
  \node[oui]  at (1*\W,-6*\H) {chaque décision};
  \node[oui]  at (2*\W,-6*\H) {chaque décision};
  \node[oui]  at (3*\W,-6*\H) {chaque décision\\(sans secrets)};

  % Légende
  \node[oui,minimum width=6mm,text width=4mm,minimum height=4mm] (l1) at (0.4,-8.75) {};
  \node[right=1mm of l1,font=\sffamily\scriptsize] (t1) {couvert};
  \node[part,minimum width=6mm,text width=4mm,minimum height=4mm,right=5mm of t1] (l2) {};
  \node[right=1mm of l2,font=\sffamily\scriptsize] (t2) {partiel};
  \node[non,minimum width=6mm,text width=4mm,minimum height=4mm,right=5mm of t2] (l3) {};
  \node[right=1mm of l3,font=\sffamily\scriptsize] {non couvert --- dit tel quel dans la raison de la décision};
\end{tikzpicture}

\caption{Ce que l'outil voit : chaque couche et ce qu'elle couvre, selon l'endroit où la commande agit}
\label{fig:vision}
\end{figure}

Trois lectures de cette figure méritent d'être soulignées.
\begin{itemize}
  \item \textbf{Seule la lecture voit partout} --- mais seulement ce qu'elle sait lire. Tout
    le reste du dispositif existe pour ce qu'elle ne sait pas lire.
  \item \textbf{La récupération et le filet s'arrêtent au projet.} Hors du projet, aucune copie
    n'est possible, et la décision le dit ; une destruction connue hors projet est escaladée, et
    celle que fait du code non lu est seulement enregistrée par le filet Linux (V18).
  \item \textbf{Le distant ne se récupère pas, il se décide.} Pour une base ou un dépôt
    distant, \demo{} ne peut que reconnaître la surface, demander ou refuser, et indiquer où
    le fournisseur garde son propre historique.
\end{itemize}
